Four sensitivity studies accompany the main grid. Their groups are small (three or four seeds), so we describe patterns rather than test them.

\begin{table}[t]\centering\caption{Training-set size: CE test IoU (mean $\pm$ std, $n$ seeds), all tasks and sizes run.}\label{tab:size}
\resizebox{\columnwidth}{!}{\begin{tabular}{llccccc}
\toprule
Method & Task & miou $\uparrow$ & iou\_circle $\uparrow$ & iou\_square $\uparrow$ & iou\_triangle $\uparrow$ & $n$ \\
\midrule
CE n=100 & none & 0.68 $\pm$ 0.04 & 0.64 $\pm$ 0.06 & 0.55 $\pm$ 0.05 & 0.56 $\pm$ 0.06 & 4 \\
CE n=250 & none & 0.76 $\pm$ 0.03 & 0.77 $\pm$ 0.04 & 0.66 $\pm$ 0.04 & 0.65 $\pm$ 0.05 & 3 \\
CE n=2000 & none & \textbf{0.82 $\pm$ 0.03} & \textbf{0.84 $\pm$ 0.02} & \textbf{0.74 $\pm$ 0.03} & \textbf{0.71 $\pm$ 0.05} & 4 \\
CE n=1000 & none & \underline{0.81 $\pm$ 0.03} & \underline{0.83 $\pm$ 0.03} & \underline{0.72 $\pm$ 0.03} & \underline{0.70 $\pm$ 0.05} & 3 \\
\midrule
CE n=100 & boundary\_p0.3 & \underline{0.67 $\pm$ 0.04} & \underline{0.63 $\pm$ 0.05} & \underline{0.54 $\pm$ 0.05} & \underline{0.55 $\pm$ 0.05} & 4 \\
CE n=2000 & boundary\_p0.3 & \textbf{0.82 $\pm$ 0.01} & \textbf{0.84 $\pm$ 0.02} & \textbf{0.73 $\pm$ 0.02} & \textbf{0.71 $\pm$ 0.02} & 4 \\
\midrule
CE n=100 & flip\_p0.3 & 0.49 $\pm$ 0.03 & 0.35 $\pm$ 0.07 & \underline{0.30 $\pm$ 0.03} & 0.32 $\pm$ 0.03 & 4 \\
CE n=250 & flip\_p0.3 & \underline{0.54 $\pm$ 0.02} & 0.47 $\pm$ 0.04 & \textbf{0.31 $\pm$ 0.02} & \underline{0.39 $\pm$ 0.03} & 3 \\
CE n=2000 & flip\_p0.3 & \textbf{0.56 $\pm$ 0.11} & \textbf{0.54 $\pm$ 0.11} & 0.30 $\pm$ 0.11 & \textbf{0.41 $\pm$ 0.24} & 4 \\
CE n=1000 & flip\_p0.3 & 0.53 $\pm$ 0.14 & \underline{0.53 $\pm$ 0.11} & 0.23 $\pm$ 0.20 & 0.36 $\pm$ 0.27 & 3 \\
\bottomrule
\end{tabular}
}\end{table}

\textbf{Training-set size (Table~\ref{tab:size}, Fig.~\ref{fig:size}).} Under clean and boundary-noise labels the curve is the same; under flips mIoU is nearly flat in $N$ ($\rhval{analysis/analysis/flip_p0.3/ce_miou_flip_p03_n100/mean:2}$ at $N=100$, $\rhval{analysis/analysis/flip_p0.3/ce_miou_flip_p03_n2000/mean:2}$ at $N=2000$), so the gap to clean labels widens with data (Sec.~\ref{sec:results}, H3). As a direct memorisation indicator, we log how often the final CE network agrees with the \emph{noisy} versus the \emph{clean} training labels on its own training pixels under flips. At $N=100$ it agrees more with the noisy labels (\rhval{analysis/analysis/flip_p0.3/fit_noisy_flip_n100/mean:2} vs \rhval{analysis/analysis/flip_p0.3/fit_clean_flip_n100/mean:2}), i.e. it has fitted flipped objects; at $N=500$ (\rhval{analysis/analysis/flip_p0.3/fit_noisy_flip_n500/mean:2} vs \rhval{analysis/analysis/flip_p0.3/fit_clean_flip_n500/mean:2}) and $N=2000$ (\rhval{analysis/analysis/flip_p0.3/fit_noisy_flip_n2000/mean:2} vs \rhval{analysis/analysis/flip_p0.3/fit_clean_flip_n2000/mean:2}) it agrees slightly more with the clean labels. Because every size gets the same number of optimiser steps, small $N$ means more epochs, which is a confound we did not separate from $N$ itself.

\begin{table}[t]\centering\caption{Noise level, CE, flip noise, $N=500$ (three seeds; clean and $p=0.3$ over the same three seeds).}\label{tab:noise}
\resizebox{\columnwidth}{!}{\begin{tabular}{llccccc}
\toprule
Method & Task & miou $\uparrow$ & iou\_circle $\uparrow$ & iou\_square $\uparrow$ & iou\_triangle $\uparrow$ & $n$ \\
\midrule
CE & flip\_p0.1 & \textbf{0.69 $\pm$ 0.02} & \textbf{0.69 $\pm$ 0.03} & \textbf{0.53 $\pm$ 0.04} & \textbf{0.57 $\pm$ 0.02} & 3 \\
\midrule
CE & flip\_p0.2 & \textbf{0.63 $\pm$ 0.04} & \textbf{0.62 $\pm$ 0.05} & \textbf{0.43 $\pm$ 0.06} & \textbf{0.50 $\pm$ 0.03} & 3 \\
\bottomrule
\end{tabular}
}\end{table}

\textbf{Flip-noise level (Table~\ref{tab:noise}).} Over seeds 0--2, mean mIoU is \rhval{analysis/analysis/flip_p0.3/lvl_clean_s012/mean:2} (clean), \rhval{analysis/analysis/flip_p0.3/lvl_flip_01/mean:2} ($p=0.1$), \rhval{analysis/analysis/flip_p0.3/lvl_flip_02/mean:2} ($p=0.2$) and \rhval{analysis/analysis/flip_p0.3/lvl_flip_03/mean:2} ($p=0.3$), with stds of \rhval{analysis/analysis/flip_p0.3/lvl_flip_sd_01/mean:2}, \rhval{analysis/analysis/flip_p0.3/lvl_flip_sd_02/mean:2} and \rhval{analysis/analysis/flip_p0.3/lvl_flip_sd_03/mean:2} for the three noisy levels: the loss is smooth and monotone in the mean across the three levels run, but the spread grows sharply at $p=0.3$. Only these levels were run, so we make no claim about other $p$, and no boundary-noise level sweep was run. The \texttt{main} group also has band-ignore CE at $p=0.1,0.2$ (Table~\ref{tab:main}); there it again tracks lower than CE at the same level.

\begin{table}[t]\centering\caption{GCE exponent $q$ under flip $p=0.3$, $N=500$ (three seeds). $q=0.7$ is the four-seed main row.}\label{tab:q}
\resizebox{\columnwidth}{!}{\begin{tabular}{lccccc}
\toprule
Method & miou $\uparrow$ & iou\_circle $\uparrow$ & iou\_square $\uparrow$ & iou\_triangle $\uparrow$ & $n$ \\
\midrule
GCE q=0.3 & \textbf{0.54 $\pm$ 0.11} & \textbf{0.51 $\pm$ 0.14} & \textbf{0.36 $\pm$ 0.08} & \textbf{0.32 $\pm$ 0.25} & 3 \\
GCE q=0.9 & \underline{0.45 $\pm$ 0.15} & \underline{0.48 $\pm$ 0.14} & \underline{0.23 $\pm$ 0.25} & \underline{0.12 $\pm$ 0.21} & 3 \\
\bottomrule
\end{tabular}
}\end{table}

\textbf{GCE exponent (Table~\ref{tab:q}).} mIoU is \rhval{sweep_gce_q/gce-q-0.3/flip_p0.3/miou/mean:2} at $q=0.3$ and \rhval{sweep_gce_q/gce-q-0.9/flip_p0.3/miou/mean:2} at $q=0.9$, against \rhval{main/gce/flip_p0.3/miou/mean:2} for the default $q=0.7$ (different seed set). The largest $q$ is clearly worst in mean (triangle IoU \rhval{sweep_gce_q/gce-q-0.9/flip_p0.3/iou_triangle/mean:2}, close to a collapse), consistent with the known under-fitting of MAE-like losses; $q=0.3$ is close to CE (\rhval{main/ce/flip_p0.3/miou/mean:2}), as expected since GCE tends to CE as $q\to0$. The main-table GCE therefore uses a default setting that was not chosen on this sweep, and the best $q$ is not identified.

\begin{table}[t]\centering\caption{CE training length at $N=100$ under flip $p=0.3$ (three seeds; the 600-step row is in Table~\ref{tab:size}).}\label{tab:steps}
\resizebox{\columnwidth}{!}{\begin{tabular}{lccccc}
\toprule
Method & miou $\uparrow$ & iou\_circle $\uparrow$ & iou\_square $\uparrow$ & iou\_triangle $\uparrow$ & $n$ \\
\midrule
CE steps=200 & \underline{0.34 $\pm$ 0.01} & \underline{0.23 $\pm$ 0.11} & \underline{0.22 $\pm$ 0.09} & \underline{0.00 $\pm$ 0.00} & 3 \\
CE steps=1200 & \textbf{0.49 $\pm$ 0.02} & \textbf{0.30 $\pm$ 0.05} & \textbf{0.36 $\pm$ 0.03} & \textbf{0.32 $\pm$ 0.04} & 3 \\
\bottomrule
\end{tabular}
}\end{table}

\textbf{Training length (Table~\ref{tab:steps}).} At $N=100$, 200 steps give mIoU \rhval{sweep_steps/ce-steps-200/flip_p0.3/miou/mean:2} (triangle IoU \rhval{sweep_steps/ce-steps-200/flip_p0.3/iou_triangle/mean:2}: the class is not yet learned), 600 steps give \rhval{analysis/analysis/flip_p0.3/ce_miou_flip_p03_n100/mean:2}, and 1200 steps give \rhval{sweep_steps/ce-steps-1200/flip_p0.3/miou/mean:2}. So mIoU rises from 200 to 600 steps and is essentially unchanged from 600 to 1200 steps: we did \emph{not} observe the late-training collapse that the early-learning literature predicts~\cite{arpit2017closer,liu2020early}. The network may begin to memorise flips only later (the train-fit statistics above show it does fit them at $N=100$), but the mIoU at 1200 steps does not show a loss.
